% !TeX program = pdflatex
% arXiv submission draft.  Suggested classes: primary math.AT; secondary cs.RO, cs.DC.
\documentclass[12pt]{article}

\setlength{\oddsidemargin}{.5cm}
\setlength{\evensidemargin}{.5cm}
\setlength{\textwidth}{6in}
\setlength{\textheight}{8in}
\setlength{\emergencystretch}{1.5em}

\usepackage{amsmath,amssymb,amsthm}
\usepackage{array,longtable}
\usepackage{float}
\usepackage{algorithm}
\usepackage{algpseudocode}
\usepackage[font=small,labelfont=bf]{caption}
\usepackage{booktabs}
\usepackage{hyperref}
\hypersetup{colorlinks=true,linkcolor=blue,citecolor=blue,urlcolor=blue}

\newcommand{\doi}[1]{doi:\,\href{http://dx.doi.org/#1}{#1}}
\newcommand{\TC}{\mathrm{TC}}
\newcommand{\SC}{\mathrm{SC}}
\newcommand{\SCst}{\SC_{\mathrm{strict}}}
\newcommand{\SD}{\mathrm{SD}}
\newcommand{\gen}[1]{\langle #1\rangle}
\newcommand{\Zt}{\mathbb{Z}_2}

\newtheorem{theo}{Theorem}[section]
\newtheorem{lemma}[theo]{Lemma}
\newtheorem{corollary}[theo]{Corollary}
\newtheorem{proposition}[theo]{Proposition}
\newtheorem{definition}[theo]{Definition}
\newtheorem{remark}[theo]{Remark}
\newtheorem{conjecture}[theo]{Conjecture}

\title{Explicit covers and rigid obstructions for simplicial complexity:\\
 from spheres to surfaces}

\author{M. Perez\footnote{Escuela Superior de F\'isica
y Matem\'aticas del Instituto Polit\'ecnico Nacional, Edificio 9, Unidad
Profesional Adolfo L\'opez Mateos, 07300 Mexico City, Mexico.}, and
J. Flores\footnotemark[1]}

\date{}

\begin{document}
\maketitle

\begin{abstract}
Gonz\'alez's strict simplicial complexity $\SCst(K)$ at subdivision level
zero is one less than the minimum number of contiguity domains covering the
ordered product $K\times K$.  Computing it needs explicit covers, and
certifying it needs lower bounds that hold for contiguity chains of every
length.  We develop both sides.  \emph{Upper bounds.}  The star cover
(every product of two vertex stars is a contiguity domain) gives
$\SCst(K)\le\rho(K)-1$ for an exactly computable set-cover constant $\rho(K)$
in seconds, sharp on every even sphere $\partial\Delta^{n+1}$.  A
chain-first GPU annealer that evolves whole systems of $k$ chains, seeded by
the star cover and descended in $k$, finds three-domain covers of
$S^2,\dots,S^6$ and of the cube boundary and descends from $24$ to $5$
domains on the icosahedron boundary; an exact SAT decision procedure for
``is this facet set a domain'' then completes the annealer's $k=4$ stall
into a verified $4$-domain cover, $\SCst(\text{icosahedron})\le3$.
\emph{Lower bounds.}  A subcomplex of $K\times K$ is \emph{rigid} if the two
projections restricted to it are joined by no contiguity chain; $k$ domains
cannot cover $K\times K$ unless $k$ classes of facets avoid every rigid
subcomplex, a finite colouring problem.  Rigid hollow triangles live over
essential empty triangles of $K$, and we prove that on any complex with no
dominated vertex the graph of every order-preserving simplicial self-map is
rigid, which contains the rigid spheres of the simplex boundaries.  The
resulting chain-length-free bounds give $\SCst(\partial\Delta^{n+1})=2>1=\TC(S^n)$
for $n=3,5,7,9$, $\SCst=2>1$ for the five-vertex M\"obius band, $\SCst\ge3>2$
for the seven-vertex torus, $4\le\SCst\le5$ for the hemi-icosahedron and
$4\le\SCst\le6$ for the hemicube.  They also explain the annealer's stalls:
on these surfaces and on $\partial\Delta^4$ every uncovered facet of a
stalled state closes a rigid subcomplex in every domain, and every star
rectangle of the hemi-icosahedron, the M\"obius band and the torus is a
maximal domain; on the icosahedron no such obstruction exists and the stall
was a search failure.  Every cover reported is verified link by link by an
auditor sharing no code with the search.
\end{abstract}

\medskip\noindent
\emph{Keywords and phrases:} topological complexity; simplicial
complexity; contiguity chains; piecewise linear motion planning;
simulated annealing; SAT; GPU computing.

\medskip\noindent
\emph{2020 Mathematics Subject Classification:} Primary: 55M30, 68W20.
Secondary: 68T40, 55U10, 90C27.

\section{Introduction}\label{sec:intro}

Farber's topological complexity $\TC(X)$ \cite{F03} measures the minimum
number of continuous motion-planning rules on a space $X$; we use the
reduced convention, so a cover by $k$ domains certifies $\TC\le k-1$.
Gonz\'alez's simplicial complexity \cite{G18} discretises it: on the ordered
product $K\times K$ of a triangulation $K$ of $X$, continuity becomes
contiguity of simplicial maps (Section~\ref{sec:back}).  The resulting
invariant $\SCst(K)$ is finite and checkable, satisfies
$\SCst(K)\ge\TC(\|K\|)$, and equals $\TC$ in the limit of barycentric
subdivisions.  At subdivision level $0$ it is an invariant of the
triangulation, and almost nothing was known about it beyond the circle and
the wedge of two circles \cite{OLGB21}.

This paper computes or brackets $\SCst$ at level $0$ for the simplex
boundaries up to $S^9$, three triangulations of $S^2$, the vertex-minimal
M\"obius band, the $7$-vertex torus, the $8$-vertex Klein bottle and four
triangulations of $\mathbb{RP}^2$.  The method has two halves, and the
paper's organisation follows them.

\emph{Upper bounds are explicit covers.}  Section~\ref{sec:star} proves that
every product $D(v,w)$ of two vertex stars is a contiguity domain
(Theorem~\ref{thm:prod}), so the exact minimum $\rho(K)$ of a small set-cover
instance gives $\SCst(K)\le\rho(K)-1$ without search; on
$\partial\Delta^{n+1}$ this is $\SCst\le2$ for all $n$, sharp for even $n$.
Section~\ref{sec:annealer} describes the chain-first annealer, which anneals
whole systems of $k$ contiguity chains on one GPU under an energy that grades
every uncovered facet by the defect count of its best chain; it finds
three-domain covers of $S^2$ through $S^6$ and of the cube boundary, where
$\rho=8$.  Section~\ref{sec:descent} documents the star-seeded descents in
$k$ on the surfaces.  Section~\ref{sec:exact} adds exact tools: a SAT
decision ``does this facet set admit a chain of at most $L$ links'', which
completes stalled annealer states, and a lazy cover solver built on it.

\emph{Lower bounds are rigid subcomplexes.}  Section~\ref{sec:rigid}
develops the notion from the companion note \cite{PM26b}: a subcomplex $Z$
of $K\times K$ is rigid if $\pi_1|_Z$ and $\pi_2|_Z$ are joined by no
contiguity chain of any length.  No domain contains a rigid subcomplex, so if
$k$ classes of facets cannot avoid a family of rigid subcomplexes then $k$
domains cannot cover (Theorem~\ref{thm:bound}).  Two sources supply the
families: hollow triangles over essential empty triangles of $K$, and, new
here, the graphs $\Gamma_w$ of order-preserving simplicial self-maps
$w\ne\mathrm{id}$ of a complex without dominated vertices
(Theorem~\ref{thm:graphs}).  The second source contains the rigid spheres
that prove $\SCst(\partial\Delta^{n+1})=2$ for $n=3,5,7,9$ and shows that the
odd spheres and the surfaces share one combinatorial property, minimality in
the sense of Barmak and Minian \cite{BM12}.  Section~\ref{sec:stall} uses
both families to dissect every stalled state of the annealer: the stalls on
$\mathbb{RP}^2$, the hemicube, the Klein bottle and $\partial\Delta^4$ are
rigid-closed dead ends, and the star rectangles of three surfaces are
maximal domains (Theorem~\ref{thm:maximal}); the stall on the icosahedron was
not an obstruction, and its exact completion gives $\SCst\le3$.

Three kinds of statement appear and are never conflated: \emph{audited
certificates} (explicit covers, re-checked by an independent program),
\emph{computer-assisted theorems} (finite SAT instances built from proved
rigid families, or witness lists re-verified from the definitions), and
\emph{empirical observations} (plateaus and stalls).  Section~\ref{sec:results}
collects the values, Section~\ref{sec:repro} the certificates and tools, and
Section~\ref{sec:concl} the open questions.

\section{Background and notation}\label{sec:back}

Throughout, $K$ is a finite, nonempty, connected simplicial complex on an
ordered vertex set $V$; a \emph{facet} is a maximal simplex, and
$\operatorname{star}(v)$ is the set of facets containing $v$.  The
\emph{ordered product} $K\times K$ is the complex of all nonempty chains
$S\subseteq V\times V$ (in the product order) with $\pi_1(S)\in K$ and
$\pi_2(S)\in K$; its realisation is the staircase triangulation of
$\|K\|\times\|K\|$, and $\pi_1,\pi_2$ are simplicial.  Its facets are the
monotone lattice paths in $s\times t$ for pairs $(s,t)$ of facets of $K$;
when $K$ is pure of dimension $d$ there are $\binom{2d}{d}$ of them over
every \emph{cell} $(s,t)$.  A set $S\subseteq V\times V$ is a simplex of
$K\times K$ iff it is a chain with both projections in $K$; hence two
simplices lie in a common facet iff their union is a simplex.

Two simplicial maps $\varphi,\varphi':J\to K$ are \emph{contiguous},
$\varphi\sim\varphi'$, if $\varphi(\sigma)\cup\varphi'(\sigma)\in K$ for
every facet $\sigma$ of $J$; a \emph{contiguity chain} of $L$ links is a
sequence $\varphi=\varphi_0\sim\cdots\sim\varphi_L=\varphi'$.  A subcomplex
$D\subseteq K\times K$ is a \emph{contiguity domain} if $\pi_1|_D$ and
$\pi_2|_D$ are joined by a contiguity chain of simplicial maps $D\to K$, and
\[
\SCst(K)=\bigl(\text{least number of domains covering }K\times K\bigr)-1
\;\ge\;\TC(\|K\|)
\]
\cite{G18,OLGB21}, with equality after enough barycentric subdivisions.  A
cover by $k$ domains is a system of $k$ piecewise-linear motion planners.
All statements here are at subdivision level $0$, and coverage figures count
facets of $K\times K$.

Two elementary facts are used repeatedly.  \emph{Heredity}: a subcomplex of
a domain is a domain (restrict the chain).  \emph{Facet classes}: if
$D_1,\dots,D_k$ cover $K\times K$ and $\mathcal F_j$ is the set of facets in
$D_j$, then $\gen{\mathcal F_j}\subseteq D_j$ is a domain and the
$\mathcal F_j$ cover the facets; so the existence of a $k$-cover is a
question about $k$ sets of facets, which may overlap but can be made
disjoint.

\section{Upper bounds without search: the star cover}\label{sec:star}

Write $c_x$ for the constant map to the vertex $x$.

\begin{theo}[product domains]\label{thm:prod}
For vertices $v,w$ of $K$ let $D(v,w)\subseteq K\times K$ be generated by
the facets over cells $(s,t)$ with $v\in s$, $w\in t$.  Then $D(v,w)$ is a
contiguity domain: if $v=v_0,v_1,\dots,v_m=w$ is an edge path of $K$, then
$\pi_1\sim c_{v_0}\sim c_{v_1}\sim\cdots\sim c_{v_m}\sim\pi_2$ on $D(v,w)$.
\end{theo}

\begin{proof}
A facet $\sigma$ over $(s,t)$ has $\pi_1(\sigma)=s$ and $\pi_2(\sigma)=t$
(a staircase path visits every row and column) and $c_x(\sigma)=\{x\}$, so
all maps are simplicial.  Consecutively, $s\cup\{v\}=s$ since $v\in s$,
$\{v_i,v_{i+1}\}$ is an edge, and $\{w\}\cup t=t$.
\end{proof}

\begin{corollary}\label{cor:rho}
Let $\rho(K)$ be the least size of a set $R\subseteq V\times V$ such that
every cell $(s,t)$ has some $(v,w)\in R$ with $v\in s$, $w\in t$.  Then
$\SCst(K)\le\rho(K)-1$.  For every $n\ge1$, $\rho(\partial\Delta^{n+1})=3$;
hence $\SCst(\partial\Delta^{n+1})\le2$, with equality $=\TC(S^n)=2$ for
even $n\ge2$.
\end{corollary}

\begin{proof}
The first claim is Theorem~\ref{thm:prod}.  On $\partial\Delta^{n+1}$ with
facets $s_a=V\setminus\{a\}$, the pairs $(0,0),(1,1),(2,2)$ cover: $(s_a,s_b)$
is missed by $(v,v)$ iff $v\in\{a,b\}$, and $\{0,1,2\}\not\subseteq\{a,b\}$.
Two pairs $(v,w),(v',w')$ never suffice: $(s_v,s_{w'})$ is missed by both.
Equality for even $n$ is Farber's $\TC(S^n)=2$ \cite{F03} and $\SCst\ge\TC$.
\end{proof}

\begin{lemma}[counting bounds]\label{lem:rows}
Let $\tau$ be the least size of a set of vertices meeting every facet,
$d=\max_v|\operatorname{star}(v)|$, and $F$ the number of facets.  Then
$\lceil\tau F/d\rceil\le\rho(K)\le\tau^2$.
\end{lemma}

\begin{proof}
A transversal $S$ of size $\tau$ gives the cover $S\times S$.  For the lower
bound fix an optimal $R$ and a facet $s$; $W_s=\{w:\exists v\in s,(v,w)\in
R\}$ is a transversal, so $|W_s|\ge\tau$, and counting incidences
$\sum_{(v,w)\in R}|\operatorname{star}(v)|\ge\sum_s|W_s|\ge\tau F$ with each
summand at most $d$.
\end{proof}

The lemma computes $\rho$ outright for the simplex boundaries, the
hemi-icosahedron ($\rho=6$) and the icosahedron ($\tau=6$, $d=5$, $F=20$,
so $\rho\ge24$, attained).  Table~\ref{tab:rho} lists $\rho$ for the
complexes of this paper; every value is the exact optimum of the set-cover
instance (CP-SAT \cite{ORTools}, status OPTIMAL, at most $12$\,s).  Exactness matters: greedy
needs $27$ rectangles on the icosahedron against $24$.

\begin{table}[H]
\centering\footnotesize
\setlength{\tabcolsep}{4pt}
\begin{tabular}{@{}lrrrrrrr@{}}
\toprule
complex $K$ & $|V|$ & $F$ & $\tau$ & $d$ &
$\lceil\tau F/d\rceil$ & $\rho(K)$ & $\tau^2$ \\
\midrule
$\partial\Delta^{n+1}$, $n\ge1$ & $n{+}2$ & $n{+}2$ & $2$ & $n{+}1$ & $3$ & $3$ & $4$\\
wedge $S^1\vee S^1$ & $4$ & $5$ & $2$ & $3$ & $4$ & $4$ & $4$\\
M\"obius band $K_M$ & $5$ & $5$ & $2$ & $3$ & $4$ & $4$ & $4$\\
cube boundary & $8$ & $12$ & $3$ & $6$ & $6$ & $8$ & $9$\\
hemi-icosahedron ($\mathbb{RP}^2$) & $6$ & $10$ & $3$ & $5$ & $6$ & $6$ & $9$\\
hemicube ($\mathbb{RP}^2$) & $7$ & $12$ & $3$ & $6$ & $6$ & $7$ & $9$\\
hemi-octahedron ($\mathbb{RP}^2$) & $9$ & $16$ & $3$ & $6$ & $8$ & $9$ & $9$\\
hemi-dodecahedron ($\mathbb{RP}^2$) & $10$ & $18$ & $4$ & $9$ & $8$ & $15$ & $16$\\
Cs\'asz\'ar torus & $7$ & $14$ & $4$ & $6$ & $10$ & $11$ & $16$\\
Klein bottle (Cervone) & $8$ & $16$ & $4$ & $7$ & $10$ & $12$ & $16$\\
icosahedron boundary & $12$ & $20$ & $6$ & $5$ & $24$ & $24$ & $36$\\
dodecahedron boundary & $20$ & $36$ & $6$ & $8$ & $27$ & $36$ & $36$\\
\bottomrule
\end{tabular}
\caption{The rectangle constant $\rho(K)$ and the bounds of
Lemma~\ref{lem:rows}; facet lists in \texttt{results/*\_starcover*/summary.json}.
The hemicube and hemi-octahedron are subdivided to be simplicial (a centre
vertex per square, a midpoint per edge cell); the pentagons of the
hemi-dodecahedron are fanned.  Three triangulations of $S^2$ carry
$\rho=3,8,36$: $\rho$ is combinatorial, not a homotopy invariant.}
\label{tab:rho}
\end{table}

\emph{The star cover as a method} (\texttt{src/star\_cover.py},
\texttt{src/rect\_cover.py}).  Solve the set cover exactly, chain each
rectangle along a shortest edge path, assign every facet to the first
rectangle containing it, shorten each chain by the greedy reduction of
\cite[Alg.~2]{OLGB21}, verify.  It cannot fail and costs seconds (the
icosahedron in $1.3$\,s greedy, $12$\,s exact; $\partial\Delta^9$ with
$1\,287\,000$ facets by three domains); it returns a complete verified cover,
the bound $\rho-1$, and chains that seed the annealer.  The gap between
$\rho-1$ and $\SCst$ is the restriction to product-shaped domains with
constant chains; it is large on sparse complexes (cube $7$ against $2$,
icosahedron $23$ against at most $3$) and closes exactly where a lower bound
rises to meet it (even spheres).

\section{Upper bounds by search: the chain-first annealer}\label{sec:annealer}

\paragraph{Why not grow domains one facet at a time.}  The pipeline of
\cite{OLGB21} grows one domain at a time by budgeted random walks
(\textsc{LocalSearch}).  Our faithful GPU port of it on
$\partial\Delta^3\times\partial\Delta^3$ ($96$ facets, walk budget $40000$,
$16$ replicas) froze every replica at a \emph{near-rectangle} of $30$ to
$54$ facets, and $65$ consecutive retry passes over a frozen front produced
zero successes in $82\%$ of the wall time; an earlier C++/MPI version on $96$
cores never completed one cover in $120$-hour jobs (authors' runs, not
measurements of \cite{OLGB21}).  Two mechanisms are consistent with the
record: an exhausted extension returns only failure, never a graded
distance to success, while contiguity is a hard constraint; and sequential,
irrevocable growth leaves fronts at which no tested extension succeeded.  In
retrospect, Theorem~\ref{thm:maximal} below shows that on several surfaces
a rectangle-shaped front \emph{cannot} be extended by any facet at all.

\paragraph{The annealer.}  A \emph{state} is a tuple of $k$ chains
$C=(\varphi_0,\dots,\varphi_{c_{\max}})$ of vertex assignments
$V\times V\to V$ with $\varphi_0=\pi_1$, $\varphi_{c_{\max}}=\pi_2$; the
interior entries are arbitrary.  A facet $\sigma$ is \emph{covered} by $C$
when every link is contiguous on $\sigma$; on the subcomplex
$\mathrm{Dom}(C)$ generated by the covered facets the entries are simplicial
and consecutive entries contiguous (each image of a covered facet lies in the
union of images of an adjacent link, a simplex of $K$), so $C|_{\mathrm{Dom}(C)}$
is a contiguity chain and $\mathrm{Dom}(C)$ a domain.  With $b_j(\sigma)$
the number of links of chain $j$ violating contiguity on $\sigma$,
$m(\sigma)=\min_j b_j(\sigma)$, $\mathrm{pen}(0)=0$ and
$\mathrm{pen}(b)=1+\beta b$ for $b>0$, the energy is
\[
E=\sum_\sigma w_\sigma\,\mathrm{pen}(m(\sigma))+\gamma\sum_j\sum_\sigma b_j(\sigma),
\]
with facet weights $w_\sigma=1$ unless a guided restart raises them.  The
first sum vanishes exactly at a cover and grades every uncovered facet by
its best chain, the graded signal the greedy lacks; the second penalises
defects in all chains and suppresses drift ($\gamma=0$ made long chains
fail in the sphere runs).  \emph{Coverage}, the reported quantity, is the
number of facets with $m(\sigma)=0$; the stopping criterion is full coverage
followed by verification, never $E=0$.

A population of $S$ states anneals in parallel on the GPU under a
Metropolis rule \cite{KGV83}.  States are initialised by random sweeps flipping $\pi_1$
towards $\pi_2$.  A proposal changes the value of one vertex in one chain at
a single interior position or over a block of two to eight consecutive
positions (\emph{track} moves, probability $p_{\mathrm{track}}$), copying
from a neighbouring map with probability $1/2$ and otherwise drawing a fresh
vertex; the surface campaigns add compound \emph{stretch}, \emph{collapse}
and \emph{seam} proposals that rewrite several columns of a block at once.
The temperature follows a geometric schedule within each epoch; every second
epoch the weakest quarter of the population is overwritten by the elite and a
few slots re-initialised.  The proposal kernel is not symmetric and we claim
no convergence guarantee.  When the incumbent covers every facet, facets are
partitioned (each to the first chain covering it), empty domains discarded,
chains reduced, and the system verified map by map, link by link, facet by
facet; the output is a certificate.  Algorithm~\ref{alg:anneal} gives the loop.

\begin{algorithm}[H]
\caption{\textsc{ChainAnneal}($K$, $k$, $c_{\max}$; $S$, $\beta$, $\gamma$,
$p_{\mathrm{track}}$).}
\label{alg:anneal}
\begin{algorithmic}[1]
\State initialise $S$ states by random sweeps from $\pi_1$ to $\pi_2$;
$\mathcal C^*\gets$ incumbent (highest coverage)
\For{epoch $=1,2,\dots$}
\ForAll{states, in parallel on the GPU}
\State $M$ Metropolis steps at temperature $T_0\to T_1$: pick a chain $j$
and a vertex $w$, propose a new value at one position or over a block of
$2$--$8$ positions (copied from a neighbour w.p.\ $1/2$, else fresh), accept
w.p.\ $\min\{1,e^{-\Delta E/T}\}$
\EndFor
\State update $\mathcal C^*$ by unweighted coverage
\If{$\mathcal C^*$ covers every facet}
\State partition facets to chains, discard empty domains, reduce each chain
\cite[Alg.~2]{OLGB21}, verify; \Return the certificate ($\SCst\le k-1$)
\EndIf
\State every second epoch: overwrite the weakest quarter with the elite,
re-initialise a few slots
\EndFor
\end{algorithmic}
\end{algorithm}

\paragraph{Ladders.}  A \emph{warm} rung initialises half the population
from the previous incumbent (chains reduced on their domains and padded with
$\pi_2$); a \emph{guided} rung additionally raises the weight of facets the
incumbent misses by $g$ after each epoch (capped at $g_{\max}$, decaying
towards $1$ by a factor $0.7$).  Rungs pay a fixed overhead
$\tau\approx40$\,s and their gain is front-loaded (an empirical fit on the
dodecahedron gives $\Delta\mathrm{cov}(t)\approx c\,t^{0.8}$), so coverage
per GPU-hour peaks at rungs of about $4\tau\approx2.7$\,min.  The surface
campaigns also used a \emph{$k$-descent} ladder: a complete $k$-cover seeds
the $(k{-}1)$-search by keeping its $k-1$ largest domains
(Section~\ref{sec:descent}).

\begin{table}[H]
\centering\footnotesize
\setlength{\tabcolsep}{3.5pt}
\begin{tabular}{@{}llr@{\hspace{8pt}}crl>{\raggedright\arraybackslash}p{3.1cm}@{}}
\toprule
target & $K$ & facets & $k$ & coverage & wall time & certified domains \\
\midrule
$S^1$ & $\partial\Delta^2$ & $18$ & $2$ & $18/18$ & $7.9$\,s & $9+9$; $\SCst=1=\TC$\\
$S^2$ & $\partial\Delta^3$ & $96$ & $3$ & $96/96$ & $22$\,s & $54+29+13$; $\SCst=2=\TC$\\
$S^3$ & $\partial\Delta^4$ & $500$ & $3$ & $500/500$ & $153$\,s & $309+130+61$; $\SCst\le2$\\
$S^4$ & $\partial\Delta^5$ & $2520$ & $3$ & $2520/2520$ & $13$\,min & $1634+653+233$; $\SCst=2=\TC$\\
$S^5$ & $\partial\Delta^6$ & $12348$ & $3$ & $12348/12348$ & $34$\,min & $7061+4239+1048$; $\SCst\le2$\\
$S^6$ & $\partial\Delta^7$ & $59136$ & $3$ & $59136/59136$ & ${\approx}2.5$\,h$^\S$ & $44909+11451+2776$; $\SCst=2=\TC$\\
\midrule
$S^1\vee S^1$ & wedge & $50$ & $3$ & $50/50$ & $41$\,s & $18+16+16$; $\SCst=2$\\
$S^2$ & cube boundary & $864$ & $3$ & $864/864$ & ${\approx}24$\,min$^\dagger$ & $482+236+146$; $\SCst=2=\TC$\\
M\"obius & $K_M$ & $150$ & $4$ & $150/150$ & $17$\,s$^\ddagger$ & $54+37+34+25$; $\SCst\le3$\\
\midrule
$S^2$ & $\partial\Delta^3$ & $96$ & $2$ & $84/96$ & plateau & two-rectangle value\\
$S^3$ & $\partial\Delta^4$ & $500$ & $2$ & $460/500$ & plateau & two-rectangle value ($8$ runs)\\
$S^4$ & $\partial\Delta^5$ & $2520$ & $2$ & $2380/2520$ & $12$\,min & two-rectangle value\\
\bottomrule
\end{tabular}
\caption{Chain-first annealer on one Tesla T4 unless noted; $k$ chains
export at most $k$ nonempty domains.  Parameters: $c_{\max}=20$--$40$,
$S=4096$ (circle $1024$), $5000$--$20000$ moves per epoch, $\beta=5$,
$\gamma=0.5$, $p_{\mathrm{track}}=0.3$ (M\"obius: $\beta=1$,
$p_{\mathrm{track}}=0.5$).  $^\dagger$Three rungs: cold $303$\,s
($662/864$), warm $814$\,s ($820$), guided $304$\,s (complete); a guided
cold rung alone stalled at $858$.  $^\ddagger$Warm from a $k=5$ cover
($36$\,s); eight $k=3$ runs stalled at $148/150$ (Section~\ref{sec:stall}).
$^\S$Eight-rung ladder on an RTX PRO 6000 ($S=64$ cold, $S=256$ guided).  The
$k=2$ plateaus stop at the coverage of the best two rectangles; the domains
realising them are often not rectangles.}
\label{tab:annealer}
\end{table}

\paragraph{The dodecahedron.}  On the dodecahedron boundary ($20$
vertices, $36$ triangles, $7776$ product facets, simplex density $3.2\%$,
dual diameter $7$) a $20$-hour campaign at $k=3$ (five warm escalations of
up to an hour, then $172$ rungs of $2.5$\,min over $12$ GPU-hours) rose from
$644$ to $2602$ facets ($33.5\%$) with mean gain per rung falling from $11.7$
to below $1$.  No cover below $\rho=36$ is known; the star bound $\SCst\le35$
stands.

\section{Star-seeded descents on the surfaces}\label{sec:descent}

The surface campaigns ran \texttt{annealer\_stretch.py} ($\beta=5$,
$\gamma=0.1$, $c_{\max}=40$, $S=4096$, $5000$ steps per epoch, seed $8233$
unless stated) with the $k$-descent ladder.  Table~\ref{tab:descent} records
first covers, descents and first stalls.  The cold first cover was the stage
that did not scale ($2$\,h on the Klein bottle; cold runs on the
hemi-octahedron and hemi-dodecahedron ended without a cover); the star cover
replaces it: the hemi-dodecahedron star-seeded at $k=14$ covered in
$6689$\,s where cold runs at $k=15,20,30$ had not after $3521$, $14673$ and
$2863$\,s, and the icosahedron descended from the $24$-domain star cover to a
verified $5$-domain cover in $17.9$\,h of rungs, every rung from $k=23$ to
$k=5$ closing (the $k=5$ rung after a warm continuation of a $2395/2400$
stall).  Single seeds, one lineage per complex: these are recorded
observations, not controlled comparisons.

\begin{table}[H]
\centering\footnotesize
\setlength{\tabcolsep}{3pt}
\begin{tabular}{@{}l r>{\raggedright\arraybackslash}p{2.9cm}>{\raggedright\arraybackslash}p{2.5cm}>{\raggedright\arraybackslash}p{1.9cm}>{\raggedright\arraybackslash}p{2.4cm}@{}}
\toprule
complex ($|K{\times}K|$) & $\rho$ & first cover & descent (wall) & lowest verified & first stall \\
\midrule
hemi-icosahedron ($600$) & $6$ & cold $k{=}11$: $679$\,s & $k{=}11\to6$ ($0.8$\,h) & $6$ domains & $k{=}5$: $540$--$544/600$\\
hemicube ($864$) & $7$ & cold $k{=}20$: $2470$\,s & $k{=}20\to7$ ($2.0$\,h) & $7$ & $k{=}6$: $836/864$\\
Klein bottle ($1536$) & $12$ & cold $k{=}15$: $7075$\,s & $k{=}15\to10$ ($4.3$\,h) & $10$ & $k{=}9$: $1494/1536$\\
hemi-octahedron ($1536$) & $9$ & star-seeded $k{=}8$: $6298$\,s & --- & $8$ & $k{=}7$: $1432/1536$ (log)\\
hemi-dodecahedron ($1944$) & $15$ & star-seeded $k{=}14$: $6689$\,s & $k{=}14\to7$ ($12.3$\,h) & $7$ & $k{=}6$: $1922/1944$\\
icosahedron ($2400$) & $24$ & star-seeded $k{=}23$: $1545$\,s & $k{=}23\to5$ ($17.9$\,h) & $5$; then $4$ (Section~\ref{sec:exact}) & $k{=}4$: $2376/2400$ ($3$ runs)\\
dodecahedron ($7776$) & $36$ & none below $\rho$ & --- & star only & $k{=}3$: $2602/7776$\\
Cs\'asz\'ar torus ($1176$) & $11$ & none below $\rho$ & --- & star only & $k{=}11$ cold: $563/1176$\\
\bottomrule
\end{tabular}
\caption{First covers, descents and first stalls.  Sources:
\texttt{results/rp2\_k11}\dots\texttt{rp2\_k6},
\texttt{hemicube\_k20}\dots\texttt{k7}, \texttt{klein\_k15\_8233\_v2}\dots\texttt{klein\_k10},
\texttt{hemi\_octa\_k8}, \texttt{hemi\_dodeca\_k14}\dots\texttt{k7},
\texttt{icosa\_k23}\dots\texttt{icosa\_k5\_v2}; stalls in \texttt{rp2\_k5\_*},
\texttt{hemicube\_k6\_v2}, \texttt{klein\_k9\_v4}, \texttt{hemi\_dodeca\_k6\_v2},
\texttt{icosa\_k4\_v3}, \texttt{dodeca\_*}.  A stall is the best coverage of
the recorded runs at that $k$; Section~\ref{sec:stall} decides which stalls
are obstructions.}
\label{tab:descent}
\end{table}

\section{Lower bounds for chains of every length: rigid subcomplexes}\label{sec:rigid}

A failed search, and even a solver infeasibility for chains of at most $L$
links, says nothing about longer chains.  The bounds of this section have no
such defect.  Proofs omitted here are in the companion note \cite{PM26b}.

\begin{definition}\label{def:rigid}
A subcomplex $Z\subseteq K\times K$ is \emph{rigid} if $\pi_1|_Z$ and
$\pi_2|_Z$ are joined by no contiguity chain of simplicial maps $Z\to K$, of
any length; equivalently, $Z$ is not a contiguity domain.  For a family
$\mathcal Z$ of subcomplexes, a \emph{$\mathcal Z$-avoiding $k$-cover} is a
family $\mathcal F_1,\dots,\mathcal F_k$ of sets of facets of $K\times K$
covering all facets such that no $\gen{\mathcal F_j}$ contains a member of
$\mathcal Z$; $\chi(\mathcal Z)$ is the least such $k$.
\end{definition}

\begin{theo}[rigid-subcomplex bound]\label{thm:bound}
If $\mathcal Z$ is a family of rigid subcomplexes then $\SCst(K)+1\ge\chi(\mathcal Z)$.
Deciding whether a $\mathcal Z$-avoiding $k$-cover exists is a finite SAT
instance (variables ``facet $f$ in class $j$'' and ``simplex $z$ absent from
class $j$'').  For $r\ge|V|^2$ the family of all rigid subcomplexes with at
most $r$ vertices gives $\chi=\SCst(K)+1$ exactly.
\end{theo}

\begin{proof}
The facet classes of a $k$-cover (Section~\ref{sec:back}) generate domains,
and by heredity a domain contains no rigid subcomplex.  Conversely, if no
class contains any rigid subcomplex, each class is itself a domain.
\end{proof}

\subsection{Source 1: hollow triangles with mismatched winding}

For a $\Zt$-valued $1$-cocycle $u$ of $K$ and a closed walk $f:C_m\to K$,
let $w_u(f)=\sum_i u(f(y_i)f(y_{i+1}))$.  Contiguous maps $C_m\to K$ have the
same winding (add the cocycle relation on the two triangles $a,b,b'$ and
$a,a',b'$ of each contiguity square $\{a,b,a',b'\}$ and sum over the cycle).
Hence a hollow triangle $Z\subseteq K\times K$ (three pairwise comparable
grid points whose pairs are simplices and whose triple is not) is rigid as
soon as $w_u(\pi_1\circ Z)\ne w_u(\pi_2\circ Z)$; the basic case is the
\emph{slice} $\tau\times\{c\}$ over an \emph{essential empty triangle}
$\tau$ (three edges of $K$, not a face, with $u(\partial\tau)=1$), where
$\pi_1$ winds and $\pi_2$ is constant.  More generally, a closed walk
$Z$ in the $1$-skeleton of $K\times K$ is rigid whenever its two projected
walks lie in different \emph{contiguity classes} of maps $C_m\to K$; we
compute these classes exactly (union--find over all $|V|^m$ maps).  They
are finer than homotopy: the $3$-walks of the Klein bottle fall into $11$
classes but its $4$-walks into $6$, the torus separates the two integral
windings ($7$ classes of $3$-walks), the hemi-octahedron has $25$ classes of
$3$-walks and $2$ of $4$-walks, while the hemi-icosahedron has exactly $2$
classes (the $\Zt$ winding) for every length $3$ to $6$ and the icosahedron
and cube boundaries have a single class for every length up to $5$ and $6$:
no rigid short cycle exists there.

\subsection{Source 2: rigid graphs over minimal complexes}

A vertex $v$ is \emph{dominated} if some $v'\ne v$ lies in every facet
containing $v$; $K$ has no dominated vertex iff
$\bigcap\{\sigma\text{ facet}:v\in\sigma\}=\{v\}$ for all $v$.  Such
complexes are the \emph{minimal} complexes of Barmak--Minian \cite{BM12}.
Every closed surface triangulation has no dominated vertex (the facets at $v$
are the triangles around the link cycle of $v$, of length $\ge3$), nor has
$\partial\Delta^{n+1}$ nor the five-vertex M\"obius band; more generally a
complex without free faces has none, since deleting a dominated vertex is a
strong collapse.

\begin{lemma}\label{lem:idalone}
If $K$ has no dominated vertex and $\varphi:K\to K$ is simplicial and
contiguous to $\mathrm{id}_K$, then $\varphi=\mathrm{id}_K$.
\end{lemma}

\begin{proof}
For a facet $\sigma$, $\sigma\cup\varphi(\sigma)\in K$ and $\sigma$ is
maximal, so $\varphi(\sigma)\subseteq\sigma$; hence
$\varphi(v)\in\bigcap_{\sigma\ni v}\sigma=\{v\}$.
\end{proof}

\begin{theo}[rigid graphs]\label{thm:graphs}
Let $K$ have no dominated vertex and let $w:V\to V$ be non-decreasing for
the vertex order and simplicial ($w(\sigma)\in K$ for all $\sigma\in K$),
with $w\ne\mathrm{id}$.  Then
\[
\Gamma_w=\bigl\{\{(v,w(v)):v\in\sigma\}:\sigma\in K\bigr\}
\]
is a subcomplex of $K\times K$, $\pi_1|_{\Gamma_w}$ is an isomorphism onto
$K$ under which $\pi_2|_{\Gamma_w}$ becomes $w$, and $\Gamma_w$ is rigid.
The same holds for the transpose $\Gamma_w^{\mathsf T}=\{\{(w(v),v):v\in\sigma\}\}$.
\end{theo}

\begin{proof}
For $\sigma\in K$, $S_\sigma=\{(v,w(v)):v\in\sigma\}$ is a chain because $w$
is non-decreasing, with $\pi_1(S_\sigma)=\sigma\in K$ and
$\pi_2(S_\sigma)=w(\sigma)\in K$ because $w$ is simplicial; so $S_\sigma$ is a
simplex of $K\times K$, and $S_\tau\subseteq S_\sigma$ for $\tau\subseteq\sigma$.
The bijection $(v,w(v))\mapsto v$ carries $S_\sigma$ to $\sigma$, hence is an
isomorphism $\Gamma_w\to K$ under which $\pi_1$ becomes $\mathrm{id}_K$ and
$\pi_2$ becomes $w$.  A contiguity chain
$\pi_1|_{\Gamma_w}=\psi_0\sim\cdots\sim\psi_L=\pi_2|_{\Gamma_w}$ therefore
becomes a chain $\mathrm{id}=\varphi_0\sim\varphi_1\sim\cdots\sim\varphi_L=w$
of simplicial maps $K\to K$; by Lemma~\ref{lem:idalone} and induction every
$\varphi_i=\mathrm{id}$, so $w=\mathrm{id}$, a contradiction.  For the
transpose read the chain backwards.
\end{proof}

\begin{remark}\label{rem:graphs}
(a) For $K=\partial\Delta^{n+1}$ every non-decreasing $w\ne\mathrm{id}$ is
simplicial (it is not onto, so proper subsets go to proper subsets) and
$\Gamma_w$ is the hollow sphere $Z_w$ of \cite{PM26b}; Theorem~\ref{thm:graphs}
extends those rigid spheres to every closed surface.  (b) Constants $w\equiv c$
give $\Gamma_c=K\times\{c\}$: no domain contains a full row or column of
$K\times K$.  (c) A contractible complex without free faces (a dunce-hat
triangulation) has $K\times\{c\}$ rigid although $\pi_1$ and $\pi_2$ are
homotopic there, so $\SCst(K)\ge1>0=\TC$: the rigidity is combinatorial.  On
the complexes of this paper, however, every $\Gamma_w$ has even mod-$2$
degree or acts non-trivially on $H^1(K;\Zt)$, so it is also homotopically
rigid (checked: no exception among $3891$ maps on the icosahedron, $597$ on
the Klein bottle, $181$ on the hemi-icosahedron).  (d) We verified by brute
force that $\mathrm{id}$ is alone in its contiguity class on all eleven
complexes of Table~\ref{tab:rho}, while a cone or a path has $27$ or $4$
maps contiguous to the identity.
\end{remark}

\subsection{The bounds}

For each complex we built the family $\mathcal Z$ of all rigid closed walks
of length $\le4$ (classes computed exactly) together with all $\Gamma_w$ and
transposes, and decided the avoidance instances with CaDiCaL
(\texttt{scripts/rigid\_analysis/rigid.py}).  Table~\ref{tab:floors} gives
the results; an UNSAT row is a theorem for chains of every length.

\begin{table}[H]
\centering\footnotesize
\setlength{\tabcolsep}{4pt}
\setlength{\tabcolsep}{3pt}
\begin{tabular}{@{}lrrrrcc@{}}
\toprule
$K$ & rigid walks & rigid graphs & $\TC{+}1$ & floor & UNSAT at $k$ & SAT at $k$ \\
\midrule
M\"obius band $K_M$ & $4\,684$ & $150$ & $2$ & $\mathbf 3$ & $2$ & $3$ \\
$\partial\Delta^4$ ($S^3$) & $0$ & $250$ & $2$ & $\mathbf 3$ & $2$ & $3$ \\
hemi-icosahedron & $19\,214$ & $362$ & $4$ & $\mathbf 5$ & $4$ ($28$\,s) & $5$ ($540$\,s) \\
hemicube & $29\,216$ & $900$ & $4$ & $\mathbf 5$ & $4$ ($486$\,s) & $5$ ($411$\,s) \\
Klein bottle & $94\,540$ & $1\,194$ & $5$ & $5$ & $4$ ($911$\,s) & $5$: undecided \\
Cs\'asz\'ar torus & $87\,204$ & $686$ & $3$ & $\mathbf 4$ & $3$ & $4$ \\
hemi-octahedron & $30\,922$ & $2\,056$ & $4$ & $4$ & $3$ & $4$ ($1476$\,s) \\
hemi-dodecahedron & $37\,966$ & $2\,042$ & $4$ & $4$ & $3$ & $4$ \\
cube, icosahedron & $0$ & $844$, $7\,782$ & $3$ & $2$ & --- & $2$ \\
\bottomrule
\end{tabular}
\caption{Chain-length-free floors from Theorem~\ref{thm:bound} with the
families of Sources 1 and 2 (distinct edge sets of rigid walks of length
$\le4$; rigid graphs with transposes); ``floor'' is in domains, and the best
known covers are in Table~\ref{tab:surfaces}.  Bold floors exceed $\TC+1$.  The
family proves nothing on the two flag triangulations of $S^2$.  For
$\partial\Delta^{n+1}$ the note \cite{PM26b} decides the rigid-sphere family
at $k=2$: UNSAT for $n=3,5,7,9$ (six, six, three and two solvers agreeing).}
\label{tab:floors}
\end{table}

\begin{theo}[computer-assisted]\label{thm:floors}
At subdivision level $0$:
$\SCst(\partial\Delta^{n+1})=2>1=\TC(S^n)$ for $n=3,5,7,9$;
$\SCst(K_M)=2>1=\TC$ for the five-vertex M\"obius band (a hand proof is in
\cite{PM26b});
$3\le\SCst(\text{Cs\'asz\'ar torus})\le10$ with $\TC=2$;
$4\le\SCst(\text{hemi-icosahedron})\le5$ and $4\le\SCst(\text{hemicube})\le6$
with $\TC(\mathbb{RP}^2)=3$ \cite{FTY03};
$4\le\SCst(\text{Klein})\le9$ with $\TC=4$ \cite{CV17}.
\end{theo}

\begin{proof}
Lower bounds: Table~\ref{tab:floors} and Theorem~\ref{thm:bound}, the
families being rigid by the winding argument and Theorem~\ref{thm:graphs}.
Upper bounds: Corollary~\ref{cor:rho} for the spheres, the certificates of
Tables~\ref{tab:annealer}, \ref{tab:descent} and~\ref{tab:surfaces}, and the
star cover for the torus.
\end{proof}

The first statement upgrades the solver result of Section~\ref{sec:sat}
(no two-domain cover of $\partial\Delta^4$ with chains of at most $48$
links) to an unconditional one, and the M\"obius value, previously
conditional on chains of at most $8$ links \cite{PM26}, is now
unconditional.  For the hemi-icosahedron the full $H^1$-cycle relaxation
(every class free of winding cycles of any length,
\texttt{scripts/relaxation/h1cegar.py}) is also satisfiable at $k=5$, so no
winding argument can exclude a five-domain cover; whether one exists is open
(Section~\ref{sec:exact}).

\section{Anatomy of a stall}\label{sec:stall}

The families of Section~\ref{sec:rigid} decide whether a stalled state can
be extended.  For a state with domains $D_1,\dots,D_k$ and uncovered facets
$L$, call $f\in L$ \emph{dead} if $D_j\cup\{f\}$ contains a rigid subcomplex
for every $j$: then no chain of any length covers $D_j\cup\{f\}$, and the
state cannot gain $f$ without first losing facets.

\begin{theo}[maximal rectangles; computer-verified]\label{thm:maximal}
Let $K$ be the hemi-icosahedron, the five-vertex M\"obius band or the
Cs\'asz\'ar torus.  For every vertex pair $(v,w)$ and every facet
$f\notin D(v,w)$, the subcomplex generated by $D(v,w)\cup\{f\}$ contains a
hollow triangle whose two projected $3$-walks lie in different contiguity
classes of maps $C_3\to K$.  Hence every star rectangle $D(v,w)$ is a
\emph{maximal} contiguity domain: no domain of any cover properly contains it.
\end{theo}

\begin{proof}
$36\cdot450$, $25\cdot96$ and $49\cdot960$ explicit witnesses
(\texttt{scripts/rigid\_analysis/}\allowbreak\texttt{certificates/}\allowbreak\texttt{maximal\_rectangles\_*.json}),
each re-verified independently: every edge of the triangle lies in a facet of
$D(v,w)\cup\{f\}$, and the two walks are in different classes, the classes
recomputed by brute force from the definition of contiguity ($2$, $3$ and $7$
classes).  Each witness is a hand proof by the winding argument: on the
hemi-icosahedron, for $D(1,1)$ and the facet over $(\{0,4,5\},\{2,3,5\})$,
the points $(0,0),(0,2),(5,5)$ span a hollow triangle whose $\pi_1$-walk
$0,0,5$ has winding $0$ and whose $\pi_2$-walk $0,2,5$ is the essential empty
triangle $\{0,2,5\}$.  Rigidity passes to supersets by heredity.
\end{proof}

On the hemicube $16$ of $49$ rectangles are maximal, on the Klein bottle $4$
of $64$.  Theorem~\ref{thm:maximal} is the \emph{product-family trap}: the
rectangle chains $\pi_1\sim c_v\sim\pi_2$ carry no defect penalty, so the
annealer is drawn to the rectangle family, and once there no single-chain
move of any width can gain a facet.  The hemi-icosahedron plateau at
$k=\rho-1=5$ ($540/600$: five diagonal rectangles, the $60$ private facets of
the sixth missing) is therefore a strict local maximum of coverage under all
single-chain moves.

\begin{table}[H]
\centering\footnotesize
\setlength{\tabcolsep}{3.5pt}
\begin{tabular}{@{}>{\raggedright\arraybackslash}p{4.3cm}crrr>{\raggedright\arraybackslash}p{2.2cm}>{\raggedright\arraybackslash}p{2.9cm}@{}}
\toprule
state & $k$ & covered & $|L|$ & dead & witnesses & exact completion \\
\midrule
\texttt{rp2\_k5\_hybrid} & $5$ & $540/600$ & $60$ & $60/60$ & hollow triangles & none ($L\le12$)\\
\texttt{rp2\_k5\_from\_k6\_\dots7232} & $5$ & $544/600$ & $56$ & $54/56$ & hollow triangles & none\\
\texttt{rp2\_k3\_cold} & $3$ & $398/600$ & $202$ & $200/202$ & hollow triangles & ---\\
\texttt{hemicube\_k6\_v2} & $6$ & $836/864$ & $28$ & $27/28$ & hollow triangles & none\\
\texttt{klein\_k9\_v4} & $9$ & $1494/1536$ & $42$ & $41/42$ & $3$-walks, one $4$-walk & none\\
\texttt{klein\_k8\_r1} & $8$ & $1483/1536$ & $53$ & $53/53$ & $3$-, $4$-walks & ---\\
\texttt{hemi\_dodeca\_k6\_v2} & $6$ & $1922/1944$ & $22$ & $6/22$ & $3$-walks & none (doors in four domains)\\
\texttt{mobius\_seed\_999} & $3$ & $148/150$ & $2$ & $0/2$ & --- & none\\
\texttt{mobius\_k3\_cold\_s4}, \texttt{\dots\_stretch\_from\_149} & $3$ & $149/150$ & $1$ & $0/1$ & --- & \textbf{found} ($8$ links)\\
\texttt{s3\_k2\_v2\_point}, \texttt{s3\_k2\_A}, \texttt{s3\_k2\_v2\_track} & $2$ & $460/500$ & $40$ & $40/40$ & rigid spheres & none ($L\le8$)\\
\texttt{icosa\_k4\_v3} & $4$ & $2376/2400$ & $24$ & $0/24$ & --- & \textbf{found} ($8$ links): \texttt{icosa\_k4\_exact}\\
\bottomrule
\end{tabular}
\caption{Door audit of the stalled states (\texttt{rigid.py audit}):
``dead'' counts uncovered facets whose addition to \emph{every} domain
closes a rigid subcomplex of the family of Table~\ref{tab:floors}; the full
mod-$2$ cohomological balance test of \cite{PM26b} agrees on every row where
it was run.  The last column is the exact decision of
Section~\ref{sec:exact} for $D_j\cup L$ over all $j$.}
\label{tab:audit}
\end{table}

Table~\ref{tab:audit} separates the stalls into two kinds.  On the
hemi-icosahedron, the hemicube, the Klein bottle and $\partial\Delta^4$ the
states are \emph{rigid-closed}: (almost) every door is shut by a hollow
triangle or a rigid sphere, the state is a dead end of any search, and only
a re-partition of the domains can continue.  On the icosahedron, the
hemi-dodecahedron and the M\"obius $149$-states the doors are \emph{open}:
no known obstruction forbids the extension, and the exact method of the next
section settles what the annealer could not.

\section{Exact methods}\label{sec:exact}

\paragraph{Deciding a single domain.}  Whether a facet set $D$ admits a
chain of at most $L$ links is a SAT instance: one-hot variables for the
$L-1$ interior maps on the grid points of $D$, and for every facet and link a
set of ``value present'' variables on which every minimal non-face of $K$
(non-edges, empty triangles, and for $\partial\Delta^{n+1}$ the vertex set)
is forbidden (\texttt{scripts/rigid\_analysis/domain\_chain\_sat.py}).
Padding by $\pi_2$ makes ``no chain of exactly $L$ links'' equivalent to
``none of at most $L$''.  Instances with up to $1\,000$ facets and $L=12$
solve in under a second; the encoder reproduces every recorded domain and
every rigid witness of Theorem~\ref{thm:maximal} (UNSAT for all tested $L$).

\paragraph{Completing a stall.}  For each domain $D_j$ of a stalled state we
ask whether $D_j\cup L$ is a domain.  On \texttt{icosa\_k4\_v3} the answer is
yes for the domain of $744$ facets at $L=8$ (and no for the smallest domain
at $L\le12$); the system of the three untouched chains and the new one is a
verified cover.

\begin{proposition}\label{prop:icosa}
$\SCst(\text{icosahedron boundary})\le3$, by the audited certificate
\texttt{results/icosa\_k4\_exact}: four domains of $819$, $728$, $629$ and
$224$ facets with chains of $8$, $9$, $7$, $6$ links.  Hence
$2\le\SCst(\text{icosahedron})\le3$.
\end{proposition}

The same step completes both M\"obius $149$-states into $3$-covers (the
one domain with an open door absorbs the last facet with an $8$-link chain),
and fails, as it must, on every rigid-closed state of Table~\ref{tab:audit}.
The annealer's $k=4$ icosahedron plateau, reached by three runs of $200$
epochs each, was thus one chain away from completion; the completing chain
re-routes the whole domain and is not reachable by the local proposals, each
of which is uphill.

\paragraph{A lazy exact cover solver.}  Combining the two halves gives a
solver for ``does $K\times K$ admit $k$ domains with chains of at most $L$
links'' (\texttt{scripts/rigid\_analysis/}\allowbreak\texttt{cegar\_cover.py}): an outer SAT
proposes $k$ (disjoint) facet classes avoiding the rigid families; the
single-domain decision checks each class; a non-domain class yields an unsat
core, minimised to a minimal non-domain facet set, which is learned as a
clause for every class.  SAT returns chains and writes an audited
certificate; UNSAT is a statement for chains of at most $L$ links, with the
learned sets as a re-checkable certificate.  It finds the three-domain
M\"obius cover in one second (the first proposal of the outer SAT is already
a cover) where the annealer stalled in eight runs, proves $k=2$ impossible
on $K_M$ and $\partial\Delta^4$ at once, and covers $\partial\Delta^3$ by
three domains in two seconds.  On the hemi-icosahedron at $k=5$, the
hemicube at $k=5$ and the icosahedron at $k=3$ it learns minimal non-domain
sets of $3$ to $25$ facets, smaller than any member of the hand-built
families, at several minutes per iteration (the outer SAT dominates); within
a two-hour budget each it reached $22$, $6$ and $17$ iterations with $290$,
$90$ and $153$ learned sets and no verdict, every proposed partition still
containing a non-domain class.  The learned sets are kept
(\texttt{certificates/learned\_*.json}) and seed later runs.  The solver is
strongest at finding covers; for
infeasibility its learned cuts are local, and on the cube boundary at $k=2$,
where $\TC(S^2)=2$ already forbids a cover, $83$ iterations and $498$
learned sets did not reach UNSAT in $38$ minutes.

\section{CP-SAT on \texorpdfstring{$\partial\Delta^4$}{the 3-sphere model}}\label{sec:sat}

Before the rigid-sphere argument we attacked $\partial\Delta^4$ at $k=2$
with a direct CP-SAT model (\texttt{paper-overleaf/paper/scs\_sat.py}): one
vertex-valued variable per grid vertex and interior map of each of $k$
chains, a facet covered by a chain when every link's image union lies in a
facet of $K$, every facet covered by some chain.  The encoding was validated
on $1100$ random chains against the annealer's defect counts, by
round-tripping five recorded covers, and by five solver controls in both
directions (\texttt{experiments/s3\_sat\_validation\_log.md}).  It returned
INFEASIBLE for $k=2$ at $L=4,5,8,12,16,24,32,40,48$ ($12$\,s to $1988$\,s,
\texttt{logs/sat\_ladder/}), a solver result for chains of at most $48$
links.  Theorem~\ref{thm:floors} now gives the unconditional statement, so
the \emph{compression} question, whether every domain of
$\partial\Delta^4\times\partial\Delta^4$ admits a chain of at most $48$
links, is no longer needed for $\SCst(\partial\Delta^4)=2$; we record that
every chainable domain we measured needs at most $4$ links.  The model does
not scale: for $\partial\Delta^6$ at $L=3,4$ and for $\mathbb{RP}^2$ at
$k=5$, $L=4$ it exhausted $5$--$6$\,GB without a verdict, which is why the
lazy solver of Section~\ref{sec:exact} decides one domain at a time.

\section{Results}\label{sec:results}

\begin{table}[H]
\centering\footnotesize
\setlength{\tabcolsep}{3pt}
\begin{tabular}{@{}lrrrcc>{\raggedright\arraybackslash}p{4.6cm}c@{}}
\toprule
complex & $|V|$ & $|K{\times}K|$ & $\rho$ & $\TC$ & floor & best certified $\SCst\le$ (source) & audit \\
\midrule
circle $\partial\Delta^2$ & $3$ & $18$ & $3$ & $1$ & $1$ & $1$ (\texttt{circle\_probe3}); $=\TC$ & yes\\
wedge $S^1\vee S^1$ & $4$ & $50$ & $4$ & $2$ & $2$ & $2$ (\texttt{wedge\_k3\_v2}); $=\TC$ & yes\\
$\partial\Delta^{n+1}$, $n=2,4,6$ & & & $3$ & $2$ & $2$ & $2$; $=\TC$ (Cor.~\ref{cor:rho}, certificates $n\le6$) & yes\\
$\partial\Delta^{n+1}$, $n=3,5,7,9$ & & & $3$ & $1$ & $\mathbf 2$ & $2$; $=\SCst>\TC$ (Thm.~\ref{thm:floors}) & yes ($n\le6$)\\
M\"obius band & $5$ & $150$ & $4$ & $1$ & $\mathbf 2$ & $2$ (\texttt{mobius\_k3\_sat}, \texttt{mobius\_k3\_cegar\_L8}); $=\SCst$ & yes\\
hemi-icosahedron & $6$ & $600$ & $6$ & $3$ & $\mathbf 4$ & $5$ (\texttt{rp2\_k6}) & yes\\
hemicube & $7$ & $864$ & $7$ & $3$ & $\mathbf 4$ & $6$ (\texttt{hemicube\_k7}) & internal\\
hemi-octahedron & $9$ & $1536$ & $9$ & $3$ & $3$ & $7$ (\texttt{hemi\_octa\_k8}) & internal\\
hemi-dodecahedron & $10$ & $1944$ & $15$ & $3$ & $3$ & $6$ (\texttt{hemi\_dodeca\_k7}) & internal\\
Klein bottle & $8$ & $1536$ & $12$ & $4$ & $4$ & $9$ (\texttt{klein\_k10}) & yes\\
Cs\'asz\'ar torus & $7$ & $1176$ & $11$ & $2$ & $\mathbf 3$ & $10$ (star cover) & yes\\
cube boundary & $8$ & $864$ & $8$ & $2$ & $2$ & $2$ (\texttt{cube\_k3\_guided}); $=\TC$ & yes\\
icosahedron boundary & $12$ & $2400$ & $24$ & $2$ & $2$ & $3$ (\texttt{icosa\_k4\_exact}) & yes\\
dodecahedron boundary & $20$ & $7776$ & $36$ & $2$ & $2$ & $35$ (star cover) & yes\\
\bottomrule
\end{tabular}
\caption{State of knowledge at level $0$.  ``floor'' is the best lower
bound on $\SCst$ (bold where it exceeds $\TC$: Theorem~\ref{thm:floors});
$\TC$ is reduced \cite{F03,FTY03,CV17}.  ``audit'': re-checked by
\texttt{src/independent\_check.py} for this paper; ``internal'': the auditor
cannot yet rebuild the complex (facet counts $864$, $1536$ and $1944$
collide with or are absent from its table), so only the search's own
verifier was run.}
\label{tab:surfaces}
\end{table}

Four remarks.  \emph{Sharp values.}  $\SCst=\TC$ on the circle, the
wedge, the even spheres up to $S^6$ and the cube; $\SCst=\TC+1$ on the
odd spheres $S^3,S^5,S^7,S^9$ and on the M\"obius band, the level-$0$ gap
being genuine and chain-length-free.  \emph{Three triangulations of
$S^2$} behave differently: $\partial\Delta^3$ and the cube are pinned at
$2$, the icosahedron is bracketed $2\le\SCst\le3$, and $\rho$ runs
$3,8,24$.  \emph{Prediction became quantitative}: rectangle arithmetic
predicted the two-chain plateaus $460=23\cdot20$ on $\partial\Delta^4$ and
$2380=34\cdot70$ on $\partial\Delta^5$, both observed, and
Theorem~\ref{thm:maximal} now explains why rectangle-shaped fronts freeze.
\emph{Where the method is blind.}  On the hemi-octahedron and the
hemi-dodecahedron the rigid families add nothing to $\TC$ and the best
covers ($8$ and $7$ domains) are far above the floor ($4$); on the torus
no cover below $\rho=11$ is known although the floor is $4$.

\section{Certificates, auditing, reproduction}\label{sec:repro}

Every cover reported is discharged as data: chains are replayed
deterministically, reduced by \cite[Alg.~2]{OLGB21}, and asserted to connect
the two projections through pairwise contiguous simplicial maps over every
facet of every domain (\texttt{cover.txt}, \texttt{planners.txt},
\texttt{summary.json}).  An independent auditor
(\texttt{src/independent\_check.py}), a pure-Python program that rebuilds
each triangulation from its facet list and shares no code with the search,
re-runs every check and exits nonzero on any discrepancy.  It verifies
$59$ archived covers: the star covers of all complexes of
Table~\ref{tab:rho} except the three subdivided models of $\mathbb{RP}^2$
(up to $\partial\Delta^9$, $1\,287\,000$ facets), the whole icosahedron
ladder (the $24$-domain seed, $k=23,\dots,5$ and \texttt{icosa\_k4\_exact}),
the hemi-icosahedron and Klein rungs, the M\"obius certificates and the
first-phase certificates of Table~\ref{tab:annealer}
(\texttt{audit/verify\_all.py --check} in the certificates repository, which
also checks a hash manifest).  The hemicube, hemi-octahedron and
hemi-dodecahedron covers pass only the internal verifier
(Table~\ref{tab:surfaces}); they are archived and hashed separately.

Computer-assisted theorems rest on three kinds of artefact: SAT instances
built from proved rigid families (\texttt{rigid.py bound}, logs in \texttt{scripts/rigid\_analysis/}\allowbreak\texttt{certificates/}\allowbreak\texttt{bound\_*.txt}; the rigid-sphere instances of \cite{PM26b}); witness lists re-verified from the definitions
(\texttt{maximal\_rectangles\_*.json}); and single-domain decisions, each
re-checkable in seconds.  The CP-SAT infeasibilities of Section~\ref{sec:sat}
are solver outputs without proof logs and are labelled as such.

Source code (annealer, star cover, rigid-subcomplex tools, exact solvers,
auditor) is at
\begin{center}\url{https://github.com/MiguelPerezI/chain-first-annealer}\end{center}
and the archived certificates at
\begin{center}\url{https://github.com/MiguelPerezI/simplicial-complexity-certificates}\end{center}
Every cover cited in this paper is in the certificates archive; the
rigid-analysis certificates, logs and stalled states are in
\texttt{opti/scripts/rigid\_analysis/}, \texttt{opti/results/} and
\texttt{opti/experiments/} of the source repository.  The versions described
here carry the tag \texttt{arxiv-v1} in both repositories.
All annealer runs except $S^6$ used one Tesla T4; wall times are single-run
observations.  The rigid-analysis and exact computations ran on a four-core
laptop CPU.

\section{Concluding remarks}\label{sec:concl}

The expensive object turned out to be the wrong object twice, and the
missing object once.  The randomised greedy hunts chains one facet at a
time; the annealer hunts whole systems and solves the classical targets in
seconds to hours; its plateaus pointed to the star cover, which makes the
first move search-free.  What was missing was a lower bound that does not
depend on chain length.  Rigid subcomplexes supply it, and they answer the
question the stalls had posed.  The annealer is weak in two precise ways: it
cannot express the global re-routing of a chain that completes an open door
(the icosahedron, the M\"obius $149$-states), and it is drawn into the
rectangle family, whose members are maximal domains on three of the
surfaces.  Where it stalls with every door rigid-closed, it is right to
stall, and the floors show how close the certified values are to what can
be proved.

The property shared by the odd spheres and the stalling surfaces is
minimality: the identity is contiguous only to itself, so every graph
$\Gamma_w$ is rigid.  What separates them from $S^2$ is that the colouring
number of their small rigid pieces exceeds $\TC+1$, while on the icosahedron
no rigid short cycle exists and the graph family is avoidable by two
classes.  Every rigid piece found is also homotopically rigid, so the
level-$0$ excess on these complexes is a resolution effect: the facets of
$K\times K$ are too coarse to separate pieces that topology also forbids.

Open at the time of writing: whether the hemi-icosahedron admits five
domains ($\SCst\in\{4,5\}$), the hemicube five or six, the icosahedron three
($\SCst\in\{2,3\}$); whether the Klein bottle needs more than five domains
(the rigid family at $k=5$ undecided); the torus between $4$ and $11$
domains; and $\SCst(\partial\Delta^{n+1})=2$ for all odd $n$, known for
$n\le9$.

\begin{thebibliography}{10}

\bibitem{BM12}
J.~A. Barmak and E.~G. Minian.
\newblock Strong homotopy types, nerves and collapses.
\newblock {\em Discrete Comput. Geom.}, 47(2):301--328, 2012.

\bibitem{CV17}
D.~C. Cohen and L.~Vandembroucq.
\newblock Topological complexity of the Klein bottle.
\newblock {\em J. Appl. Comput. Topol.}, 1(2):199--213, 2017.

\bibitem{F03}
M.~Farber.
\newblock Topological complexity of motion planning.
\newblock {\em Discrete Comput. Geom.}, 29(2):211--221, 2003.

\bibitem{FTY03}
M.~Farber, S.~Tabachnikov, and S.~Yuzvinsky.
\newblock Topological robotics: motion planning in projective spaces.
\newblock {\em Int. Math. Res. Not.}, 2003(34):1853--1870, 2003.

\bibitem{G18}
J.~Gonz\'alez.
\newblock Simplicial complexity: piecewise linear motion planning in
  robotics.
\newblock {\em New York J. Math.}, 24:279--292, 2018.

\bibitem{KGV83}
S.~Kirkpatrick, C.~D. Gelatt, Jr., and M.~P. Vecchi.
\newblock Optimization by simulated annealing.
\newblock {\em Science}, 220(4598):671--680, 1983.

\bibitem{OLGB21}
C.~Ortiz, A.~Lara, J.~Gonz\'alez, and A.~Borat.
\newblock A randomized greedy algorithm for piecewise linear motion
  planning.
\newblock {\em Mathematics}, 9(19):2358, 2021.

\bibitem{ORTools}
L.~Perron and V.~Furnon.
\newblock OR-Tools (CP-SAT solver).
\newblock Google, \url{https://developers.google.com/optimization}.

\bibitem{PM26}
M.~Perez and J.~Flores.
\newblock The strict simplicial complexity of the vertex-minimal
  M\"obius band triangulation is~$2$.
\newblock Companion note, 2026.

\bibitem{PM26b}
M.~Perez and J.~Flores.
\newblock Rigid subcomplexes and lower bounds for strict simplicial
  complexity at level~$0$.
\newblock Supplementary note, 2026.

\end{thebibliography}

\end{document}
